\chapter{Hỏi Để Tự Nhìn Lại}
\chapterintro{Không phải câu hỏi nào cũng để lấy đáp án}{Có những câu hỏi sinh ra để học sinh soi lại cách học và cách sống của mình.}

\section{Em hay bỏ cuộc ở phút nào nhất?}
\begin{cauhoibox}{Câu hỏi}
Em hay bỏ cuộc ở phút nào nhất: phút đầu, phút giữa, hay phút gần xong?
\end{cauhoibox}
\begin{taisaocoich}
Câu hỏi cụ thể nên học sinh dễ trả lời thật hơn. Nó mở ra đối thoại về thói quen học tập mà không nặng tính lên lớp.
\end{taisaocoich}
\begin{goiybien}
Hợp cho giờ chủ nhiệm, giờ kỹ năng, hoặc lúc lớp vừa trải qua một bài khó.
\end{goiybien}
\ngancach

\section{Em muốn thầy cô bớt làm điều gì để em học tốt hơn?}
\begin{cauhoibox}{Câu hỏi}
Em muốn thầy cô bớt làm điều gì để em học tốt hơn?
\end{cauhoibox}
\begin{taisaocoich}
Câu này trả lại tiếng nói cho học sinh. Khi được hỏi thật, các em thường hợp tác hơn ở phần ngược lại: thầy cô muốn các em thay đổi gì.
\end{taisaocoich}
\begin{goiybien}
Nên nói trước: trả lời tử tế, cụ thể, không mỉa mai. Có thể cho viết giấy ẩn danh.
\end{goiybien}
\ngancach

\section{Điều gì ở em tốt hơn tuần trước một chút?}
\begin{cauhoibox}{Câu hỏi}
Điều gì ở em tốt hơn tuần trước một chút?
\end{cauhoibox}
\begin{taisaocoich}
Học sinh dễ nói thật hơn khi thầy cô không đòi thay đổi lớn. Câu hỏi này tạo cảm giác tiến bộ nhỏ nhưng có thật.
\end{taisaocoich}
\begin{goiybien}
Hợp cuối tuần, cuối tháng, hoặc sau một giai đoạn lớp vừa mệt vừa áp lực.
\end{goiybien}
\ngancach

\section{Điều gì ở cách học của em nên giữ nguyên?}
\begin{cauhoibox}{Câu hỏi}
Điều gì ở cách học của em nên giữ nguyên vì nó đang giúp em tốt lên?
\end{cauhoibox}
\begin{taisaocoich}
Tự nhìn lại không chỉ để sửa. Câu hỏi này giúp học sinh nhận ra điểm mạnh nhỏ của mình, từ đó bớt cảm giác bị soi lỗi liên tục.
\end{taisaocoich}
\begin{goiybien}
Rất hợp dùng xen giữa các câu hỏi cần thay đổi để giờ tự nhìn lại không bị nặng.
\end{goiybien}
\ngancach

\section{Nếu được xin cả lớp giúp một điều để học tốt hơn, em xin điều gì?}
\begin{cauhoibox}{Câu hỏi}
Nếu được xin cả lớp giúp một điều để học tốt hơn, em xin điều gì?
\end{cauhoibox}
\begin{taisaocoich}
Câu hỏi này mở ra nhu cầu thật của học sinh và kéo tập thể vào vai trò hỗ trợ nhau. Nó rất hợp với lớp cần tăng sự đồng cảm.
\end{taisaocoich}
\begin{goiybien}
Có thể cho trả lời ẩn danh bằng giấy nếu lớp chưa quen nói thật trước tập thể.
\end{goiybien}
